\documentclass[a4paper, 12pt]{article}

\usepackage[nopatch = all]{microtype}
\usepackage{anyfontsize, libertine, bm}

\usepackage[margin = .8in]{geometry}
\usepackage{multicol}

\setlength \parindent {0pt}
\newcommand \newslide {\par\noindent\rule{\linewidth}{.5pt}\par}
\newcommand \newframe {\par\noindent\rule{\linewidth}{.5pt}\par
                       \vskip-.9\baselineskip
                           \noindent\rule{\linewidth}{.5pt}\par}

\begin{document}

\begin{multicols*}{3}
  Good morning, everyone.
  Today, I will share the calculation of band structure of 2D Kagome Lattice
  via finite element analysis.
  \newframe
  Firstly, I will introduce the Finite Element Analysis method,
  short for ``FEA''.
  \newframe
  In short, FEA is a divide and conquer approach.
  It is a general numerical method for solving PDE in two- or
  three-space variables especially for some boundary value problems.
  \newslide
  It subdivides a large system into smaller, simpler parts called
  finite elements.
  \newslide
  When applying it to solve the PDEs, one will approximate a continuous domain,
  with a collection of simple sub-domains, which is called elements.
  \newslide
  For example, given a material.
  Then one can slice it like this figure.
  \newslide
  Then, using local ``piecewise'' polynomials instead of global functions.
  \newslide
  The field is solved only at specific points or nodes,
  also shown in this figure.
  \newframe
  Next, I'll talk about the mathematical principle.
  Schr\"odinger equation acts an important part in the strong-correlating
  electron theory. Usually, we need to solve it directly at every point, but FEA
  satisfies it on average using the method of weighted residuals.
  \newslide
  To do this, we need to obtain its weak formulation.
  Since the Schr\"odinger equation is an eigenequation,
  we can sightly rearrange the Schr\"odinger equation by defining $R$,
  then our goal is to minimize the residual.
  \newslide
  Usually, or I can say ``in the Cartesian coordinate system'', the Hamiltonian
  operator $\mathcal H$ is constructed with the kinetic term and the
  potential term. We expand it, and then multiply a test function $v$,
  finally integrate over the full space.
  \newslide
  Finally, we shall apply the integration by parts, and there will be an
  boundary term, which will be vanished due to the periodicity.
  Now, we have obtained the weak formulation.
  \newframe
  OK. Now, I will talk about the methodology.
  \newframe
  We need to use a shape function to form the unknown field.
  \newslide
  The unknown field is a linear combination of the values at the nodes
  \newslide
  and the shape functions.
  The shape functions are usually linear or quadratic polynomials.
  \newslide
  Typically, it has a value of 1 at node $i$ and 0 at all other nodes,
  as shown in the figure.
  \newslide
  Now, we switch from the nodes to the element.
  Each element has its stiffness matrix and mass matrix, as defined following.
  They can be used to construct the global matrices.
  \newslide
  by summing them.
  Then, we arrive at another eigenvalue problem:
  from the eigenvalue for Hamiltonian operator to the eigenvalue to matrix.
  \newframe
  We choose the triangular element as an example now:
  the three nodes of it are $(x_1,y_1)$, $(x_2,y_2)$, and $(x_3,y_3)$.
  Then, the shape functions will become as this.
  \newslide
  We substitute it into the expression of the stiffness Matrix and
  mass matrix we mentioned before.
  \newframe
  OK. Now, let me introduce the parameters.
  We choose the Kagome lattice (in Chinese, its Lóng Mù, which the lattice will
  be displayed later).
  We will give out its lattice geometry,
  \newslide
  numerical parameters,
  \newslide
  and the uniform physical constants.
  \newframe
  Originally, people usually tend to use the tight-binding model,
  so, this table shows the differences between it and the FEA method.
  \newslide
  Also, in the lattice, we assume the potential is periodic.
  \newslide
  The numerical workflow including defining the geometry,
  \newslide
  dividing the mesh,
  \newslide
  and physically, the next three steps are used to analyze the band structure.
  \newframe
  What's more, the wavefunction in the lattice should satisfy the Bloch's
  theorem to obey the periodic condition.
  \newslide
  In FEA, this condition will be implemented as a constraint on the boundary
  nodes of the unit cell mesh.
  \newslide
  As we can see that this relation has made the wavefunction raise a phase
  factor, which can be also implemented in FEA.
  \newslide
  This makes the global matrix dependent on the crystal momentum.
  The eigenvalue problem of the matrix becomes this one.
  \newslide
  Next, we can compare the results of the two methods.
  The tight-binding method predicts a flat band and two dispersive bands.
  \newslide
  Concerning to the FEA method,
  in low-energy limit, it recovers the same structure,
  \newslide
  But also shows higher energy bands, i.e., the excited states.
  So, this method is more robust.
  \newframe
  More important, I need to classify the nodes. It can be classified into
  \newslide
  Master nodes,
  \newslide
  slave nodes,
  \newslide
  and interior nodes due to different geometry and physics conditions.
  \newslide
  Due to the Classification of the nodes, the wavefunction at each element
  should be also classified.
  Since the master nodes and the slave nodes are both on the edges,
  their wavefunction just differ by a phase factor.
  Then, the eigenequation can be rewritten in terms of $\tilde{\bm\Psi}$.
  \newframe
  Now, we arrive at the most important part.
  We need to define the domain,
  \newslide
  and then choose the unit cell for the Kagome lattice.
  This figure shows the structure of the Kagome lattice,
  at each point, there is a carbon atom,
  \newslide
  but we need to divide them into three
  categories, labeled as $A$, $B$, and $C$.
  since the unit cell is a $C_{3h}$ group.
  \newframe
  Next, we need to mesh it.
  The mesh respects the 3-fold symmetry.
  \newslide
  Local refinement near potential wells,
  \newslide
  Periodic node pairs are identified along opposite edges.
  \newslide
  In total there are about 5000 elements, and 10000 degrees of freedom.
  \newslide
  Figure (a) shows the meshing of the unit cell with an inset highlighting the
  ligament.
  \newslide
  Figure (b) shows the schematics of the boundary conditions in the FE
  simulations to determine the equivalent longitudinal, shear, and torsional
  stiffnesses of the ligament.
  \newslide
  Inside each well, the potential is defined using the step function.
  \newslide
  This piecewise constant potential is easy to integrate in FEA.
  \newframe
  Now, I will show the results of this model.
  \newslide
  The solution of the energy band including two
  dispersive bands and one flat band.
  \newslide
  which are plotted in the left figures.
  \newslide
  People usually care about the high-symmetry points.
  In the energy band figure, the path $\Gamma \to K \to M \to \Gamma$ is
  sampled with $40$ points per segment.
  \newslide
  And as the increase of the number of the elements, the flat band energy
  converges to within 1\% with the finest mesh.
  \newframe
  Finally, we will learn about the Geometric Effects in FEA.
  \newslide
  This effect is not available in the tight-binding method.
  \newslide
  We can obtain how the band gap changes with atom size,
  \newslide
  The transition from nearly-free electrons to strongly bound states,
  \newslide
  and we can remove elements from the mesh to simulate vacancies.
  \newslide
  OK. So much for today's presentation.
  \newframe
  Thanks for listening.
\end{multicols*}

\end{document}